\documentclass[colorlinks]{thesis-kando}

% -------------------------------------------------
% 1. PDF verzió kezelése (a figyelmeztetés megszüntetéséhez)
% -------------------------------------------------
\pdfminorversion=7

% -------------------------------------------------
% 2. Magyar nyelv specifikus beállítások (a babel ELŐTT)
% -------------------------------------------------
\PassOptionsToPackage{defaults=hu-min}{magyar.ldf}

% -------------------------------------------------
% 3. Alapvető kódolás és nyelvi csomag
% -------------------------------------------------
\usepackage[T1]{fontenc}
\usepackage[utf8]{inputenc}
\usepackage[magyar]{babel}

% -------------------------------------------------
% 4. Minden egyéb csomag (a te “template”-ed szerint, duplikáció nélkül)
% -------------------------------------------------
\usepackage{graphicx,amsmath,amssymb,amsthm,hulipsum,bussproofs}
\usepackage{listings,xcolor,rotating,booktabs,fancyvrb}
\usepackage[normalem]{ulem}
\usepackage{pdfpages}
\usepackage[linguistics]{forest}
\usepackage{stackengine,scalerel}
\usepackage{url}
\usepackage{xstring} % szükséges a patch-hez

% Jobb URL tördelés (a te mintád szerint)
\def\UrlBreaks{\do\/\do-\do\.\do\a\do\b\do\c\do\d\do\e\do\f\do\g\do\h\do\i\do\j\do\k\do\l\do\m\do\n\do\o\do\p\do\q\do\r\do\s\do\t\do\u\do\v\do\w\do\x\do\y\do\z\do\0\do\1\do\2\do\3\do\4\do\5\do\6\do\7\do\8\do\9}

% Formázási és matematikai beállítások (a template-ből)
\footnotestyle{rule=fourth}
\DeclareMathOperator{\tg}{tg}

\newtheorem{tetel}{Tétel}[chapter]
\newtheorem{lemma}[tetel]{Lemma}
\theoremstyle{definition}
\newtheorem{definicio}[tetel]{Definíció}
\newtheorem{feladat}[tetel]{Feladat}
\theoremstyle{remark}
\newtheorem{megjegyzes}[tetel]{Megjegyzés}
\newtheorem*{megoldas}{Megoldás}

% -------------------------------------------------
% Képek keresési útvonalai (az \includegraphics-hoz)
% -------------------------------------------------
\graphicspath{{img/}{./img/}{./}}

% -------------------------------------------------
% SAFE INCLUDE – biztos működés img/ mappával is, foreach nélkül
%   Használat: \safeinclude[<includegraphics opciók>]{<név kiterjesztés nélkül>}{<dobozszélesség>}
%   Pl.: \safeinclude[width=5cm]{logo}{5cm}
% -------------------------------------------------
\newcommand{\safeinclude}[3][]{%
	\begingroup
	\def\found{}%
	% 1) pontos fájlnév (ha valaki kiterjesztéssel hívja)
	\IfFileExists{#2}{\def\found{#2}}{}%
	% 2) kiterjesztés nélküli név → próbálgatás root-ban
	\ifx\found\empty
	\IfFileExists{#2.png}{\def\found{#2.png}}{}%
	\fi
	\ifx\found\empty
	\IfFileExists{#2.jpg}{\def\found{#2.jpg}}{}%
	\fi
	\ifx\found\empty
	\IfFileExists{#2.jpeg}{\def\found{#2.jpeg}}{}%
	\fi
	\ifx\found\empty
	\IfFileExists{#2.pdf}{\def\found{#2.pdf}}{}%
	\fi
	% 3) img/ alatti próbálgatás
	\ifx\found\empty
	\IfFileExists{img/#2}{\def\found{img/#2}}{}%
	\fi
	\ifx\found\empty
	\IfFileExists{img/#2.png}{\def\found{img/#2.png}}{}%
	\fi
	\ifx\found\empty
	\IfFileExists{img/#2.jpg}{\def\found{img/#2.jpg}}{}%
	\fi
	\ifx\found\empty
	\IfFileExists{img/#2.jpeg}{\def\found{img/#2.jpeg}}{}%
	\fi
	\ifx\found\empty
	\IfFileExists{img/#2.pdf}{\def\found{img/#2.pdf}}{}%
	\fi
	% 4) ./img/ alatti próbálgatás
	\ifx\found\empty
	\IfFileExists{./img/#2}{\def\found{./img/#2}}{}%
	\fi
	\ifx\found\empty
	\IfFileExists{./img/#2.png}{\def\found{./img/#2.png}}{}%
	\fi
	\ifx\found\empty
	\IfFileExists{./img/#2.jpg}{\def\found{./img/#2.jpg}}{}%
	\fi
	\ifx\found\empty
	\IfFileExists{./img/#2.jpeg}{\def\found{./img/#2.jpeg}}{}%
	\fi
	\ifx\found\empty
	\IfFileExists{./img/#2.pdf}{\def\found{./img/#2.pdf}}{}%
	\fi
	
	\ifx\found\empty
	\fbox{\parbox{#3}{\centering Hiányzó kép:\\ \texttt{#2}}}%
	\else
	\includegraphics[#1]{\found}%
	\fi
	\endgroup
}

% -------------------------------------------------
% PATCH: thesis-kando titlepage expects figures/Logo.jpg
%   - Javítva: a saját img/logo.* fájlt fogja használni.
% -------------------------------------------------
\makeatletter
\newcommand{\titlelogofallback}{logo} % img/logo.png vagy jpg/pdf stb.

\newcommand{\patchedmaketitle}{%
	\let\origincludegraphics\includegraphics
	\renewcommand{\includegraphics}[2][]{%
		\begingroup
		\def\argfile{##2}%
		\def\targetfile{figures/Logo.jpg}%
		\ifx\argfile\targetfile
		\safeinclude[##1]{\titlelogofallback}{8cm}%
		\else
		\origincludegraphics[##1]{##2}%
		\fi
		\endgroup
	}%
	\maketitle
	\let\includegraphics\origincludegraphics
}
\makeatother

% -------------------------------------------------
% Listings beállítása a magyar karakterekhez
% -------------------------------------------------
\lstset{
	inputencoding=utf8,
	extendedchars=true,
	literate={á}{{\'a}}1 {é}{{\'e}}1 {í}{{\'i}}1 {ó}{{\'o}}1 {ö}{{\"o}}1 {ő}{{\H{o}}}1
	{ú}{{\'u}}1 {ü}{{\"u}}1 {ű}{{\H{u}}}1
	{Á}{{\'A}}1 {É}{{\'E}}1 {Í}{{\'I}}1 {Ó}{{\'O}}1 {Ö}{{\"O}}1 {Ő}{{\H{O}}}1
	{Ú}{{\'U}}1 {Ü}{{\"U}}1 {Ű}{{\H{U}}}1,
	breaklines=true,
	basicstyle=\ttfamily\small,
	keywordstyle=\color{blue},
	commentstyle=\color{gray},
	stringstyle=\color{teal}
}

% Define JavaScript language for listings
\lstdefinelanguage{JavaScript}{
	keywords={break, case, catch, continue, debugger, default, delete, do, else, finally, for, function, if, in, instanceof, new, return, switch, this, throw, try, typeof, var, void, while, with, let, const, export, import, await, async},
	ndkeywords={class, export, boolean, throw, implements, import, this, new},
	sensitive=true,
	comment=[l]{//},
	morecomment=[s]{/*}{*/},
	morestring=[b]',
	morestring=[b]",
	alsoletter={-},
	alsodigit={:},
	morekeywords=[2]{=>, yield, async, await, constructor},
	keywordstyle=\color{blue}\bfseries,
	ndkeywordstyle=\color{green}\bfseries,
	identifierstyle=\color{black},
	commentstyle=\color{gray}\ttfamily,
	stringstyle=\color{red}\ttfamily,
}

% Customize listings appearance
\lstset{
	language=JavaScript,
	extendedchars=true,
	inputencoding=utf8,
	basicstyle=\ttfamily\small,
	numbers=left,
	numberstyle=\tiny\color{gray},
	stepnumber=1,
	numbersep=10pt,
	showstringspaces=false,
	breaklines=true,
	frame=single,
	captionpos=b,
	escapeinside={\%*}{*)},
}

% Fordított 'T' szimbólum (a template-ből)
\def\tang{\ThisStyle{\abovebaseline[0pt]{\scalebox{-1}{$\SavedStyle\perp$}}}}

% -------------------------------------------------
% Document
% -------------------------------------------------
\begin{document}
	
	% --- Címlap metaadatok ---
	\title{D\&D Ultimate Tool\\Projekt dokumentáció}
	\author{Adeniran Ádám Kristóf\\László Tamás\\Szoftverfejlesztő és -tesztelő szak}
	\institute{Szoftver fejlesztés és tesztelés}
	\supervisor{Kerényi Róbert Nándor\\oktató}
	\city{Miskolc}
	\date{2026}
	
	% Biztonságos címlap (ne álljon meg Logo.jpg hiány miatt)
	\patchedmaketitle
	
	\tableofcontents
	\newpage
	
	% =================================================
	% 1. fejezet — Bevezetés
	% =================================================
	\chapter{Bevezetés}
	
	\section{A projekt bemutatása és célja}
	A \textbf{D\&D Ultimate Tool} egy webes alkalmazás, amely támogatja a \textbf{Dungeons \& Dragons} játékosokat és mesélőket
	kampányok, karakterek és közösségi funkciók kezelésében. A rendszer célja egy egységes, többmodulos felület biztosítása
	a játék előkészítéséhez és lebonyolításához, valós idejű kommunikációval és bővíthető adatmodellel.
	
	\begin{figure}[ht!]
		\centering
		\safeinclude[width=10cm]{logo}{10cm}
		\caption[Projekt logó]{D\&D Ultimate Tool logó}
		\label{fig-logo}
	\end{figure}
	
	\section{Célközönség meghatározása}
	A célcsoport az aktív D\&D közösség:
	\begin{itemize}
		\item játékosok, akik karaktereiket és felszereléseiket strukturáltan szeretnék kezelni,
		\item mesélők (DM-ek), akik encounter-t, kampányt és jegyzeteket vezetnek,
		\item baráti társaságok, akik online szervezik a játékalkalmakat és kommunikációt.
	\end{itemize}
	
	\section{Alkalmazott technológiák és eszközök}
	Mellékletre hivatkozás:
	\newline
	Lásd: \ref{fig-melleklet1} melléklet ábrát.
	
	\begin{itemize}
		\item \textbf{Adatbázis:} MySQL
		\item \textbf{Backend:} ASP.NET Core Web API, Entity Framework
		\item \textbf{Webes kliens:} React (reszponzív kialakítás)
		\item \textbf{Fejlesztői eszközök:} Git/GitHub, VS Code, Visual Studio, Postman/Swagger
	\end{itemize}
	
	
	% =================================================
	% 2. fejezet — Technológiák (teljes, kidolgozott)
	% =================================================
	\chapter{Technológiák}
	
	\section{Használt technológiák}
	A \textbf{D\&D Ultimate Tool} fejlesztése során olyan technológiai stack-et alkalmaztunk, amely biztosítja a
	\textbf{skálázhatóságot}, a \textbf{biztonságos adatkezelést}, valamint a \textbf{gyors és modern felhasználói élményt}.
	A rendszer architektúrája három fő pillérre épül: \textbf{backend} (API és üzleti logika), \textbf{frontend}
	(felhasználói felület) és \textbf{adatbázis} (perzisztens adattárolás).
	
	A megoldás alapelve, hogy a kliensoldali alkalmazás (React) kizárólag a megjelenítésért és a felhasználói interakciókért felel,
	míg az üzleti logika és az adat-hozzáférés a szerveren (ASP.NET Core Web API) történik. Ez a szétválasztás
	(\emph{separation of concerns}) könnyebb karbantartást, jobb tesztelhetőséget és rugalmas bővíthetőséget eredményez.
	
	\begin{figure}[ht!]
		\centering
		\safeinclude[width=12cm]{architecture}{12cm}
		\caption[Rendszer áttekintése]{A rendszer áttekintése: React frontend -- ASP.NET Core API -- MySQL adatbázis}
		\label{fig-arch}
	\end{figure}
	
	\subsection{Backend}
	A backend megvalósításához \textbf{ASP.NET Core Web API}-t használunk, amely magas teljesítményű és platformfüggetlen keretrendszer.
	A backend fő felelősségei:
	\begin{itemize}
		\item \textbf{Autentikáció és jogosultságkezelés:} JWT alapú beléptetés, védett végpontok.
		\item \textbf{Üzleti logika:} karakterek, kampányok, wiki bejegyzések és dobások kezelése.
		\item \textbf{Chat és közösségi funkciók:} üzenetküldés és kapcsolati logika (barát rendszer).
		\item \textbf{Adat-hozzáférés:} entitások kezelése és lekérdezések végrehajtása.
	\end{itemize}
	
	A réteges felépítés célja, hogy a kontroller réteg csak a HTTP kérések kezelésére szolgáljon, míg a tényleges üzleti szabályok
	külön szolgáltatási (service) rétegben helyezkedjenek el. Ez csökkenti a kód ismétlődését, és megkönnyíti a tesztelést.
	
	\subsubsection{ASP.NET Core}
	Az \textbf{ASP.NET Core} moduláris felépítése lehetővé teszi a middleware alapú kérésfeldolgozást, ami jól használható
	például naplózás, CORS, hibakezelés, illetve autentikáció beállítására. A projektben különösen fontos, hogy az API-k
	konzisztens formában szolgáltassanak adatot a frontend felé (jellemzően JSON formátumban), és stabilan kezeljék a hibás kéréseket.
	
	Az alábbi példa egy tipikus végpont szerkezetét szemlélteti (szemléltető jelleggel):
	
	\begin{lstlisting}[language=JavaScript,caption={Példa API végpont (szemléltető)}]
		[HttpGet("list")]
		public async Task<IActionResult> GetFriends()
		{
			int userId = int.Parse(User.FindFirstValue(ClaimTypes.NameIdentifier)!);
			
			var friends = await _context.Friendships
			.Where(f => (f.RequesterId == userId || f.AddresseeId == userId) && f.Status == FriendshipStatus.Accepted)
			.Include(f => f.Requester)
			.Include(f => f.Addressee)
			.ToListAsync();
			
			//Itt van a helye a friend dto kezelésnek
			
			return Ok(friendDtos);
		}
		
	\end{lstlisting}
	
	\subsubsection{Entity Framework Core}
	A perzisztencia rétegben az \textbf{Entity Framework Core} (EF Core) szerepe, hogy a C\# objektumokat és az adatbázis táblákat
	összekösse, így az adatkezelés a domain modell szintjén történhet. Az EF Core támogatja a migrációkat is, ami lehetővé teszi, hogy
	a fejlesztés során az adatbázis-séma verziózottan, reprodukálható módon változzon.
	
	\subsubsection{Autentikáció (JWT)}
	A felhasználói munkamenet kezeléséhez \textbf{JWT} (JSON Web Token) alapú autentikációt alkalmazunk. A bejelentkezés során a szerver
	tokent ad vissza, amelyet a kliens a következő kérések során a \texttt{Authorization: Bearer <token>} fejlécben továbbít.
	Ez biztosítja, hogy a védett erőforrásokat csak az arra jogosult felhasználók érjék el.
	
	\subsection{Adatbázis}
	Az adatok tárolására \textbf{MySQL} relációs adatbázist alkalmazunk. A relációs modell előnye, hogy strukturált,
	konzisztens adattárolást biztosít, és jól támogatja a kapcsolatok kezelését (felhasználók--barátok, kampányok--játékosok,
	üzenetek--felhasználók, stb.). A táblák közötti kapcsolatok idegen kulcsokkal és indexekkel optimalizálhatók, ami különösen
	fontos nagyobb adatmennyiség és gyakori lekérdezések esetén.
	
	A projektben tipikus entitások:
	\begin{itemize}
		\item \textbf{Users}: felhasználói adatok
		\item \textbf{Friends}: barát kapcsolatok (állapotokkal)
		\item \textbf{Messages}: privát és csoportos üzenetek
		\item \textbf{Campaigns / Characters}: kampányok és karakterek alapadatai
	\end{itemize}
	
	A fejlesztés során az adatbázis kezeléséhez használható eszköz például \textbf{phpMyAdmin} vagy \textbf{MySQL Workbench}.
	A cél a jól átgondolt, normalizált struktúra, amely csökkenti az ismétlődést, és támogatja az adatintegritást.
	
	\subsection{Frontend}
	A felhasználói felületet \textbf{React} alapokon valósítjuk meg, amely komponens-alapú fejlesztést tesz lehetővé. A React
	előnye, hogy a UI logika újrafelhasználható elemekből épül fel, és a felület gyorsan frissíthető állapotváltozások hatására.
	A projektben a frontend tipikusan az alábbi feladatokat látja el:
	\begin{itemize}
		\item \textbf{Oldalak és nézetek:} kampányok, karakterlap, wiki, chat.
		\item \textbf{Állapotkezelés:} felhasználói session, kiválasztott kampány/karakter, UI állapotok.
		\item \textbf{API kommunikáció:} HTTP kérések kezelése, token továbbítása, hibák megjelenítése.
		\item \textbf{Reszponzív megjelenés:} mobil és asztali nézet támogatása.
	\end{itemize}
	
	\subsubsection{React}
	A React ökoszisztémája támogatja a kliensoldali routingot (pl. React Router), valamint bővíthető állapotkezelési megoldásokat
	(Context API / Redux). A komponens-alapú szerkezet segíti, hogy az alkalmazás egyes részei (pl. chat panel, karakterkártyák,
	kampánylista) külön fejleszthetők és karbantarthatók legyenek.
	
	\begin{lstlisting}[language=JavaScript,caption={Példa React komponens (szemléltető)}]
		import React from 'react';
		import { Link } from 'react-router-dom';
		import '../assets/styles/Navbar.css';
		import '../assets/styles/WikiTheme.css'; // unified theme
		
		//Itt van a helye a szekciók felsorolásának
		
		const Wiki = () => {
			return (
			<div className="page-comp">
			<div className="page-overlay">
			<h1>Wiki</h1>
			<p>Explore the different sections of the game wiki below:</p>
			
			
			//Kártya Gridek helye
			</div>
			</div>
			);
		};
		
		export default Wiki;
		
	\end{lstlisting}
	
	\subsection{Fejlesztői eszközök}
	A fejlesztés során az alábbi eszközök támogatják a hatékony munkavégzést és a minőségbiztosítást:
	\begin{itemize}
		\item \textbf{Git \& GitHub:} verziókezelés, branch stratégia, pull request alapú fejlesztés.
		\item \textbf{Visual Studio / VS Code:} backend és frontend fejlesztés, debug és refaktor eszközök.
		\item \textbf{Postman / Swagger:} API végpontok ellenőrzése és dokumentálása.
		\item \textbf{Adatbázis kliens:} phpMyAdmin / MySQL Workbench (opcionális).
	\end{itemize}
	
	% =================================================
	% 3. fejezet — Követelményanalízis
	% =================================================
	\chapter{Követelményanalízis}
	
	\section{Funkcionális követelmények}
	\begin{itemize}
		\item \textbf{Karakterkezelés:} karakterek létrehozása, módosítása, mentése
		\item \textbf{Kampánykezelés:} kampányok, kalandnapló, jegyzetek
		\item \textbf{Wiki modul:} világ és homebrew tartalom dokumentálása
		\item \textbf{DM Tools:} kezdeményezés, NPC/szörny adatbázis, encounter támogatás
		\item \textbf{Barát rendszer:} kapcsolatok kezelése
		\item \textbf{Chat:} privát üzenetek és csoportos beszélgetések
	\end{itemize}
	
	\section{Nem funkcionális követelmények}
	\begin{itemize}
		\item Reszponzív felület (mobil és desktop)
		\item Biztonságos hitelesítés és jogosultságkezelés (JWT)
		\item Skálázható, réteges architektúra
		\item Karbantartható kódbázis (elkülönített frontend/backend)
	\end{itemize}
	
	\section{Használati esetek (Use Cases)}
	Példák:
	\begin{itemize}
		\item Felhasználó regisztrál és bejelentkezik.
		\item Mesélő kampányt hoz létre, játékosokat meghív.
		\item Játékos karaktert készít, és chatben egyeztet.
	\end{itemize}
	
	% =================================================
	% 4. fejezet — Rendszertervezés
	% =================================================
	\chapter{Rendszertervezés}
	
	\section{Adatbázis terv}
	\subsection{ER-diagram és kapcsolatok}
	Az adatbázis relációs modellre épül. A barát rendszer és üzenetek kezeléséhez külön táblák tartoznak (barátkapcsolat, üzenetnapló).
	
	\subsection{Adatszótár és táblaleírások}
	A főbb entitások: Users, Friends, Messages, Campaigns, Characters, WikiEntries.
	
	\section{Rendszerarchitektúra és komponensek}
	\subsection{ORM (Object-Relational Mapping)}
	Entity Framework használata a domain modellek és a MySQL relációk leképzéséhez.
	
	\subsection{Entity Framework}
	Code First / Database First megközelítés a projekt ütemezésétől függően.
	
	\section{REST API végpontok tervezése}
	Tipikus végpontok:
	\begin{itemize}
		\item \texttt{POST /api/auth/login}
		\item \texttt{GET /api/friend/list}
		\item \texttt{POST /api/chat/send}
	\end{itemize}
	
	\section{UI/UX Tervek (Wireframes)}
	A fő képernyők: bejelentkezés, dashboard, kampány nézet, karakterlap, chat.
	
	% =================================================
	% 5. fejezet — Megvalósítás
	% =================================================
	\chapter{Megvalósítás}
	
	\section{Adatbázis réteg implementációja}
	MySQL táblák, kulcsok és indexek kialakítása, adatkonzisztencia szabályokkal.
	
	\section{REST Backend fejlesztése}
	\subsection{Adatkezelés és Üzleti logika}
	Szolgáltatási réteg (service layer) felel az üzleti szabályokért, controller réteg a végpontokért.
	
	\subsection{Biztonság és Autentikáció (JWT)}
	JWT alapú tokenkezelés, jogosultsági szintek és védett végpontok.
	
	\section{React Frontend fejlesztése}
	\subsection{Komponens-architektúra}
	Komponens alapú felépítés, külön modulok a karakter, kampány és chat területekhez.
	
	\subsection{Reszponzív megoldások és Mobil nézet}
	Grid/flex alapú elrendezés, mobilbarát menü és töréspontok.
	
	\begin{lstlisting}[language=JavaScript,caption={Példa React komponens}]
		
		{/* Card Grid */}
		<div className="wiki-grid">
		{wikiSections.map((section, index) => (
			<Link key={index} to={section.path} className="wiki-card">
			{section.name}
			</Link>
			))}
		</div>
		<h1>Detailed Information</h1>
		
		{/* Card Grid */}
		<div className="wiki-grid">
		{baseInfoSection.map((section, index) => (
			<Link key={index} to={section.path} className="wiki-card">
			{section.name}
			</Link>
			))}
		</div>
	\end{lstlisting}
	
	\section{WPF Desktop alkalmazás}
	(Amennyiben tartozik a projekthez.) MVVM struktúra, adatkezelés és kiegészítő admin modulok.
	
	% =================================================
	% 6. fejezet — Minőségbiztosítás és Tesztelés
	% =================================================
	\chapter{Minőségbiztosítás és Tesztelés}
	
	\section{Unit tesztek (Egységtesztelés)}
	Backend: xUnit vagy NUnit; kritikus üzleti logikák tesztelése.
	
	\section{API tesztelés (Postman/Swagger)}
	Swagger UI dokumentáció; Postman gyűjtemények regresszióhoz.
	
	\section{Felhasználói felület és Reszponzivitás tesztelése}
	Kézi és automatizált UI ellenőrzések (pl. Lighthouse).
	
	\section{Hibajegyzék és javítások}
	Verziózott hibakezelés GitHub Issues-ban.
	
	% =================================================
	% 7. fejezet — Felhasználói útmutató
	% =================================================
	\chapter{Felhasználói útmutató}
	
	\section{Rendszerkövetelmények és Telepítés}
	\begin{itemize}
		\item Modern böngésző (Chrome/Firefox/Edge)
		\item Backend futtatás: .NET SDK és MySQL
	\end{itemize}
	
	\section{A szoftver használata (Web/WPF)}
	\begin{enumerate}
		\item Böngészőben megnyitás
		\item Bejelentkezés
		\item Kampány/karakter létrehozás
		\item Chat és eszközök használata játék közben
	\end{enumerate}
	
	\begin{figure}[ht!]
		\centering
		\safeinclude[width=12cm]{login}{12cm}
		\caption{Bejelentkezési képernyő}
		\label{fig-login}
	\end{figure}
	
	% =================================================
	% 8. fejezet — Összegzés
	% =================================================
	\chapter{Összegzés}
	
	\section{Projekt értékelése}
	A D\&D Ultimate Tool célja egy olyan komplex rendszer biztosítása, amely egységes felületen támogatja a kampánykezelést,
	karakteradminisztrációt, dobásokat és közösségi funkciókat.
	
	\section{Továbbfejlesztési lehetőségek}
	\begin{itemize}
		\item Mobilbarát UI további fejlesztése
		\item Discord bot integráció
		\item Szabálykönyv-import és verziókövetés
		\item Többnyelvű felület támogatása
	\end{itemize}
	
	% =================================================
	% Köszönetnyilvánítás
	% =================================================
	\newpage
	\Huge\begin{center}
		\textbf{Köszönetnyilvánítás}
	\end{center}\normalsize
	Köszönet mindazoknak, akik hozzájárultak a projekt elkészítéséhez, külön megköszönve \ldots
	\par
	
	% =================================================
	% Mellékletek
	% =================================================
	\chapter*{Mellékletek}
	\addcontentsline{toc}{chapter}{Mellékletek}
	
	\begin{figure}[ht!]
		\centering
		\safeinclude[width=14cm]{rdb}{14cm}
		\caption[1. melléklet]{Adatbázis ábra}
		\label{fig-melleklet1}
	\end{figure}
	
	\clearpage
	
	\begin{figure}[p]
		\centering
		\begin{turn}{90}
			\safeinclude[width=23cm]{melleklet}{23cm}
		\end{turn}
		\caption[2. melléklet]{Szkennelt dokumentum}
		\label{fig-melleklet2}
	\end{figure}
	
	\clearpage
	
	% =================================================
	% Irodalomjegyzék
	% =================================================
	\begin{thebibliography}{9}
		\addcontentsline{toc}{chapter}{\bibname}
		
		\bibitem{Fazekas}
		\textsc{Fazekas István}: \emph{Valószínűségszámítás}, Debreceni Egyetem, Debrecen, 2004.
		
		\bibitem{Tomacs}
		\textsc{Tómács Tibor}: \emph{A valószínűségszámítás alapjai}, Líceum Kiadó, Eger, 2005.
		
		\bibitem{VartereszI1}
		\textsc{Várterész Magdolna}: \emph{Az informatika logika alapjai előadások}, 2006/07-es tanév 1.félév,
		\url{http://users.atw.hu/de-mi/index.php?r=file/download\&id=219}, Letöltve: 2019 március 12.-én
		
	\end{thebibliography}
	
\end{document}
